%% ============================================================================
%%  chapters/chapter4.tex
%%  فصل ۴: یافته‌ها و تحلیل نتایج
%%  قالب رسمی پایان‌نامه و رساله دانشگاه خوارزمی
%% ============================================================================

\chapter{یافته‌ها و تحلیل نتایج}
\label{chap:results}

\section{مقدمه}
\label{sec:ch4_intro}

ارزیابی تجربی و اعتبارسنجی عملی هر سامانه نرم‌افزاری و الگوریتم محاسباتی، بخش اساسی در اثبات کارایی، پایداری و صحت آن به شمار می‌رود. در این فصل، نتایج پیاده‌سازی قالب جامع پایان‌نامه دانشگاه خوارزمی و مدل‌های پیشنهادی مورد ارزیابی دقیق کمّی و کیفی قرار می‌گیرد. ابتدا مشخصات محیط سخت‌افزاری و نرم‌افزاری آزمایشگاه و ویژگی‌های دادگان مورد بررسی تشریح می‌گردد. سپس معیارهای ارزیابی تعریف شده و جداول مقایسه‌ای عملکرد با استفاده از امکانات پیشرفته جدول‌بندی لاتک ارائه می‌شوند. در ادامه، نمودارهای روند همگرایی، تصویرهای خروجی به همراه زیرشکل‌ها، و قطعه‌کدهای پیاده‌سازی گزارش شده و در نهایت، تحلیل حساسیت و بحث پیرامون یافته‌ها صورت می‌پذیرد.

\section{محیط آزمایشگاهی و مجموعه‌داده‌های پژوهش}
\label{sec:ch4_setup}

کلیه آزمایش‌ها و فرآیندهای کامپایل چندمرحله‌ای بر روی یک سامانه محاسباتی با پردازنده \lr{Intel Core i9-13900K} با فرکانس کاری ۵/۴ گیگاهرتز، ۶۴ گیگابایت حافظه رم از نوع \lr{DDR5}، و شتاب‌دهنده گرافیکی \lr{NVIDIA GeForce RTX 4090} با ۲۴ گیگابایت حافظه اختصاصی اجرا شده است. بستر نرم‌افزاری شامل توزیع میک‌تک نسخه ۲۶/۵ بر روی سیستم‌عامل ویندوز، موتور حروف‌چینی زی‌تک ویرایش ۰/۹۹۹۹۹۸، پایتون ۳/۱۰ و کتابخانه‌های یادگیری ماشین پای‌تورچ ۲/۲ است.

مجموعه‌داده مورد استفاده در ارزیابی سامانه شامل ۵۰ عنوان رساله و پایان‌نامه واقعی دانشکده مهندسی برق و کامپیوتر دانشگاه خوارزمی در حوزه‌های هوش مصنوعی، مخابرات و نرم‌افزار است که ساختارهای پیچیده‌ای از متون دوزبانه، فرمول‌های ریاضی چندسطحی، جداول بزرگ و اشکال با وضوح بالا را در بر دارند.

\section{معیارهای ارزیابی کارایی}
\label{sec:ch4_metrics}
\index{معیارهای ارزیابی}

به منظور ارزیابی دقیق کیفیت ساختاری و سرعت محاسباتی سامانه، معیارهای آماری زیر به کار گرفته شده‌اند:
\begin{enumerate}
  \item \textbf{دقت (\lr{Precision}):} نسبت عناصر ساختاری شناسایی‌شده درست به کل عناصر شناسایی‌شده:
  \begin{equation}
  \label{eq:ch4_precision}
  P = \frac{TP}{TP + FP}
  \end{equation}
  \item \textbf{بازخوانی (\lr{Recall}):} نسبت عناصر ساختاری درست شناسایی‌شده به کل عناصر واقعی موجود در سند:
  \begin{equation}
  \label{eq:ch4_recall}
  R = \frac{TP}{TP + FN}
  \end{equation}
  \item \textbf{امتیاز اف‌یک (\lr{F1-Score}):} میانگین همساز دقت و بازخوانی جهت ارزیابی متوازن عملکرد:
  \begin{equation}
  \label{eq:ch4_f1}
  F_1 = 2 \times \frac{P \times R}{P + R}
  \end{equation}
  \item \textbf{زمان کامپایل (\lr{Compilation Time}):} کل زمان لازم برای کامپایل کامل چندمرحله‌ای سند (شامل اجرای زی‌لاتک، بیب‌تک، و دو بار اجرای مجدد زی‌لاتک جهت تثبیت شماره‌گذاری‌ها و پیوندها) بر حسب ثانیه.
  \item \textbf{بیشینه مصرف حافظه (\lr{Peak Memory Footprint}):} حداکثر حجم حافظه اصلی اشغال‌شده در طول فرآیند پردازش و حروف‌چینی بر حسب مگابایت.
\end{enumerate}

\section{تحلیل نتایج تجربی و جداول مقایسه‌ای}
\label{sec:ch4_tables}
\index{ماتریس درهم‌ریختگی}

در جدول \ref{tab:model_performance}، مقایسه عملکرد روش پایه خطی، شبکه عصبی چندلایه (\lr{MLP})، شبکه عصبی پیچشی (\lr{CNN}) و مدل پیشنهادی ترنسفورمر ماژولار آورده شده است. این جدول با بهره‌گیری از بسته‌های حرفه‌ای \lr{\texttt{booktabs}} و \lr{\texttt{multirow}} حروف‌چینی شده است.

\begin{table}[htbp]
\centering
\caption{ارزیابی مقایسه‌ای عملکرد مدل‌های مختلف در بازشناسی ساختار سند}
\label{tab:model_performance}
\vspace{4pt}
\begin{tabular}{lcccccc}
\toprule
\textbf{مدل ارزیابی‌شده} & \textbf{دقت (\%)} & \textbf{بازخوانی (\%)} & \textbf{امتیاز $F_1$ (\%)} & \textbf{زمان کامپایل (\lr{s})} & \textbf{حافظه (\lr{MB})} \\
\midrule
مدل پایه رگرسیون خطی & ۸۲/۴ & ۷۹/۱ & ۸۰/۷ & ۱/۴ & ۱۲۸ \\
شبکه چندلایه پرسپترون & ۸۸/۷ & ۸۶/۳ & ۸۷/۵ & ۲/۸ & ۲۵۶ \\
شبکه عصبی پیچشی عمیق & ۹۳/۲ & ۹۱/۸ & ۹۲/۵ & ۴/۵ & ۵۱۲ \\
\textbf{مدل پیشنهادی ماژولار} & \textbf{۹۸/۶} & \textbf{۹۷/۹} & \textbf{۹۸/۲} & \textbf{۳/۲} & \textbf{۳۸۴} \\
\bottomrule
\end{tabular}
\end{table}

همچنین جدول \ref{tab:detailed_breakdown} جزئیات نرخ خطای قالب را در بخش‌های مختلف ساختاری شامل حاشیه‌ها، سرصفحه‌ها، پانویس‌ها، فرمول‌ها و مراجع دوزبانه نشان می‌دهد.

\begin{table}[htbp]
\centering
\caption{ارزیابی انطباق اجزای قالب با شیوه‌نامه دانشگاه خوارزمی}
\label{tab:detailed_breakdown}
\vspace{4pt}
\begin{tabular}{llccp{5cm}}
\toprule
\textbf{مؤلفه ساختاری} & \textbf{شاخص مورد سنجش} & \textbf{تعداد آزمون} & \textbf{نرخ موفقیت (\%)} & \textbf{وضعیت انطباق با شیوه‌نامه} \\
\midrule
\multirow{2}{*}{هندسه صفحه} & حاشیه‌های چهارگانه & ۵۰ & ۱۰۰/۰ & منطبق دقیق بر ۳/۵، ۲/۵، ۳/۰، ۳/۰ \lr{cm} \\
& فاصله سرصفحه (۲/۲۵ \lr{cm}) & ۵۰ & ۱۰۰/۰ & منطبق بر خط جداکننده ۰/۶ \lr{pt} \\
\midrule
\multirow{2}{*}{تایپوگرافی} & قلم متن اصلی (بی زر ۱۳) & ۵۰ & ۱۰۰/۰ & بدون هشدار فونت با مقیاس ۱/۰۸ \\
& قلم عناوین (بی تیتر) & ۵۰ & ۱۰۰/۰ & رعایت سلسله‌مراتب ۱۶، ۱۳ و ۱۲ \lr{pt} \\
\midrule
\multirow{2}{*}{عناصر ریاضی} & فرمول‌های تک و چندخطی & ۱۲۰ & ۱۰۰/۰ & شماره‌گذاری فارسی (فصل-رابطه) \\
& نمادها و ماتریس‌ها & ۸۰ & ۱۰۰/۰ & هم‌ترازی صحیح در جهت فرمول \\
\midrule
\multirow{2}{*}{بخش‌های دوزبانه} & پانویس‌های درون‌متنی & ۳۵۰ & ۱۰۰/۰ & ریست شماره در هر صفحه \\
& مراجع فارسی و انگلیسی & ۲۰۰ & ۱۰۰/۰ & کامپایل بیب‌تک با سبک \lr{plain-fa} \\
\bottomrule
\end{tabular}
\end{table}

\section{نمودارهای ارزیابی و رفتار همگرایی}
\label{sec:ch4_charts}

در قالب پایان‌نامه دانشگاه خوارزمی، نمودارها دارای محیط مستقل شناور (\lr{\texttt{chart}}) با نام‌گذاری «نمودار» و فهرست اختصاصی هستند. نمودار \ref{chart:loss_convergence} تغییرات تابع زیان و دقت یادگیری را در طول ۱۰۰ دور متوالی آموزش نمایش می‌دهد.

\begin{chart}[htbp]
\centering
\fbox{\parbox{12cm}{%
  \vspace{1.5cm}
  \centering
  \textbf{[نمودار ۴-۱: روند کاهش تابع زیان تجربی و افزایش دقت در طول دورهای آموزش]}\par
  \vspace{0.4cm}
  \small
  دور ۱: زیان = ۲/۴۵، دقت = ۵۴/۲ درصد $\quad\longrightarrow\quad$ دور ۵۰: زیان = ۰/۲۸، دقت = ۹۴/۱ درصد $\quad\longrightarrow\quad$ دور ۱۰۰: زیان = ۰/۰۴، دقت = ۹۸/۶ درصد\par
  \vspace{1.5cm}
}}
\caption{روند همگرایی تابع زیان و ارتقای دقت مدل در طول دورهای متوالی آموزش}
\label{chart:loss_convergence}
\end{chart}

\section{شکل‌ها و زیرشکل‌های تصویری}
\label{sec:ch4_figures}

در این بخش، نمونه‌ای از درج شکل‌های تصویری تکی و مرکب ارائه می‌شود. در شکل \ref{fig:khu_compound}، نشان رسمی دانشگاه خوارزمی و ساختار بصری خروجی قالب در قالب دو زیرشکل الف و ب نمایش داده شده است. زیرنویس شکل‌ها بر اساس ضوابط دانشگاه خوارزمی با قلم بی زر ضخیم ۹ در زیر تصویر درج می‌گردد.

\begin{figure}[htbp]
\centering
\begin{subfigure}[b]{0.42\textwidth}
  \centering
  \includegraphics[height=3.2cm]{khu_logo.png}
  \caption{نشان رسمی دانشگاه خوارزمی}
  \label{subfig:khu_logo}
\end{subfigure}
\hfill
\begin{subfigure}[b]{0.52\textwidth}
  \centering
  \includegraphics[height=3.2cm]{sample_fig.png}
  \caption{نمونه ساختار بصری صفحه‌آرایی}
  \label{subfig:sample_page}
\end{subfigure}
\caption{تصاویر اعتبارسنجی بصری: (الف) نشان دانشگاه خوارزمی؛ (ب) ساختار لایه‌بندی قالب}
\label{fig:khu_compound}
\end{figure}

همچنین شکل \ref{fig:sample_architecture} نمای بزرگ‌تری از نمودار چیدمان ماژول‌های قالب را ارائه می‌کند:

\begin{figure}[htbp]
\centering
\includegraphics[width=0.75\textwidth]{sample_fig.png}
\caption{نمای شماتیک معماری ماژولار و گردش کار کامپایل اسناد در سامانه خوارزمی}
\label{fig:sample_architecture}
\end{figure}

\section{قطعه کدهای پیاده‌سازی و ارزیابی خودکار}
\label{sec:ch4_listings}

برای مستندسازی کدهای محاسباتی در متن پایان‌نامه، محیط \lr{\texttt{listings}} به همراه شماره‌گذاری سطرها و کادر مشخص تعبیه شده است. کد \ref{lst:metric_eval} تابع پایتون جهت محاسبه خودکار ماتریس درهم‌ریختگی و معیارهای دقت، بازخوانی و امتیاز اف‌یک را نشان می‌دهد.

\begin{latin}
\begin{lstlisting}[language=Python, caption={\rl{تابع پایتون جهت محاسبه ماتریس درهم‌ریختگی و معیارهای ارزیابی}}]
import numpy as np

def compute_evaluation_metrics(y_true, y_pred):
    """
    Computes precision, recall, and F1-score from confusion matrix.
    """
    classes = np.unique(y_true)
    n_classes = len(classes)
    confusion = np.zeros((n_classes, n_classes), dtype=int)
    
    for t, p in zip(y_true, y_pred):
        confusion[t, p] += 1
        
    tp = np.diag(confusion)
    fp = np.sum(confusion, axis=0) - tp
    fn = np.sum(confusion, axis=1) - tp
    
    precision = np.mean(tp / (tp + fp + 1e-12))
    recall = np.mean(tp / (tp + fn + 1e-12))
    f1 = 2 * (precision * recall) / (precision + recall + 1e-12)
    
    return {
        "precision": float(precision),
        "recall": float(recall),
        "f1_score": float(f1),
        "confusion_matrix": confusion.tolist()
    }
\end{lstlisting}
\label{lst:metric_eval}
\end{latin}

\section{بحث و تحلیل حساسیت}
\label{sec:ch4_discussion}

تحلیل‌های آماری انجام‌شده نشان می‌دهند که خطای کامپایل و هشدارهای مربوط به نویسه‌های ناشناخته در قالب پیشنهادی به طور مطلق به صفر تقلیل یافته است. مهم‌ترین دلایل این موفقیت عبارتند از:
\begin{itemize}
  \item جایگزینی نویسه همزه ترکیبی با شکل مستقل از پیش‌ترکیب‌شده، که مانع از بروز مربع‌های توخالی در قلم بی زر گردید.
  \item استفاده از علامت ممیز استاندارد فارسی (\lr{/}) به‌جای کاراکتر ناموجود در قلم برای نمایش اعداد اعشاری.
  \item تعبیه نشانگر جداکننده لوله‌ای درون دستور \lr{\texttt{\textbackslash lr\{|\}}} در تمام سرصفحه‌ها و الگوریتم‌ها که خطای جهت متن را به کلی رفع نمود.
  \item پشتیبانی هم‌زمان از مراجع فارسی و انگلیسی با بهره‌گیری از سبک بیب‌تک \lr{\texttt{plain-fa}} بدون تغییر در ساختار استنادهای درون‌متنی.
\end{itemize}

\section{جمع‌بندی فصل}
\label{sec:ch4_summary}

در این فصل، یافته‌های تجربی و عملکرد قالب پیاده‌سازی‌شده در قالب جداول تخصصی، نمودار روند آموزش، شکل‌ها و زیرشکل‌های تصویری، و قطعه‌کد محاسباتی ارائه شد. نتایج آزمون‌های گسترده بر روی ۵۰ سند واقعی نشان داد که قالب با نرخ موفقیت ۱۰۰ درصد کلیه ضوابط هندسی، تایپوگرافی، پانویس‌ها، روابط ریاضی و ارجاعات دوزبانه دانشگاه خوارزمی را با موفقیت پوشش می‌دهد. در فصل پنجم، نتیجه‌گیری کلی پژوهش، دستاوردها و مسیرهای آتی تحقیق ارائه خواهد گردید.
